\section{Integritas, Autentikasi, dan Aplikasi CLI}

\subsection{Mekanisme Integritas dan Autentikasi Data}

\subsubsection{Analisis Keamanan Encrypt-then-MAC (EtM)}
Dalam kriptografi simetris, kombinasi primitif enkripsi dan kode autentikasi pesan (\textit{Message Authentication Code}/MAC) dapat diwujudkan melalui tiga paradigma utama: \textit{MAC-then-Encrypt} (MtE), \textit{Encrypt-and-MAC} (E\&M), dan \textit{Encrypt-then-MAC} (EtM).

Paradigma MtE (sebagaimana diterapkan pada SSL 3.0 / TLS 1.0) menghitung tag MAC atas plaintext, melakukan penambahan \textit{padding}, lalu mengenkripsi keduanya. Struktur ini rentan terhadap serangan \textit{Padding Oracle} (seperti serangan Vaudenay), di mana penyerang memanipulasi ciphertext dan memanfaatkan perbedaan respons server terhadap kesalahan \textit{padding} versus kesalahan MAC untuk mendekripsi ciphertext byte demi byte tanpa mengetahui kunci. Paradigma E\&M menghitung tag MAC langsung dari plaintext secara paralel dengan enkripsi, yang berpotensi membocorkan kesamaan pesan melalui tag MAC yang deterministik.

Sistem pada Tugas D secara ketat mengimplementasikan komposisi \textbf{Encrypt-then-MAC (EtM)}. Plaintext terlebih dahulu dienkripsi menggunakan kunci enkripsi $K_{\text{enc}}$, menghasilkan ciphertext $C$. Selanjutnya, tag autentikasi $T$ dihitung atas keseluruhan header dan ciphertext menggunakan kunci MAC terpisah $K_{\text{mac}}$:
\[
C = \text{Enc}_{K_{\text{enc}}}(P), \quad T = \text{MAC}_{K_{\text{mac}}}(\text{Header} \,\|\, C)
\]
Pada proses dekripsi, verifikasi tag $T$ dieksekusi sebelum mendekripsi ciphertext atau memproses \textit{padding}. Apabila ciphertext atau header mengalami modifikasi sekecil satu bit pun, tag CMAC dinyatakan tidak valid dan proses langsung dihentikan (\textit{fail-stop}). Dekriptor tidak pernah memproses data korup, sehingga \textit{Padding Oracle Attack} maupun serangan \textit{chosen-ciphertext} adaptif (IND-CCA2) dicegah secara mutlak.

\subsubsection{Konstruksi CMAC (OMAC1 / RFC 4493)}
Fungsi integritas menggunakan Cipher-based MAC (CMAC / OMAC1) berbasis block cipher 128-bit sesuai spesifikasi RFC 4493. CMAC mengatasi kelemahan CBC-MAC klasik terhadap serangan ekstensi panjang (\textit{length extension attack}) pada pesan dengan panjang variabel melalui penggunaan dua \textit{subkey} turunan, $K_1$ dan $K_2$.

Pembangkitan subkey dilakukan atas medan hingga Galois $GF(2^{128})$ dengan polinomial reduksi primitif:
\[
P(x) = x^{128} + x^7 + x^2 + x + 1 \implies R_{128} = \text{0x00}^{\dots 14} \,\|\, \text{0x87}
\]
Blok konstanta nol $0^{128}$ dienkripsi untuk menghasilkan nilai dasar $L = \text{CIPH}_{K_{\text{mac}}}(0^{128})$. Operasi perkalian dengan $x$ pada $GF(2^{128})$ diimplementasikan melalui fungsi pergeseran bit kiri satu posisi (\textit{doubling}) dengan reduksi kondisional:
\[
K_1 = \begin{cases}
L \ll 1, & \text{jika } \text{MSB}_1(L) = 0 \\
(L \ll 1) \oplus R_{128}, & \text{jika } \text{MSB}_1(L) = 1
\end{cases}
\]
\[
K_2 = \begin{cases}
K_1 \ll 1, & \text{jika } \text{MSB}_1(K_1) = 0 \\
(K_1 \ll 1) \oplus R_{128}, & \text{jika } \text{MSB}_1(K_1) = 1
\end{cases}
\]

Pesan yang diautentikasi $M$ dibagi menjadi blok-blok berukuran 16 byte $M_1, M_2, \dots, M_n$. Blok terakhir $M_n$ diproses sebagai berikut:
\begin{itemize}
    \item Jika blok terakhir lengkap ($|M_n| = 16$ byte), blok tersebut di-XOR dengan subkey $K_1$: $M_n^* = M_n \oplus K_1$.
    \item Jika blok terakhir tidak lengkap ($|M_n| < 16$ byte) atau pesan berukuran 0 byte, diterapkan \textit{padding} bit $1$ diikuti bit $0$ hingga genap 16 byte ($10^*$-padding), kemudian di-XOR dengan subkey $K_2$: $M_n^* = \text{pad}(M_n) \oplus K_2$.
\end{itemize}
Rantai CBC dieksekusi dengan IV awal bernilai $0^{128}$, dan hasil enkripsi blok terakhir menghasilkan tag CMAC berukuran 16 byte (128 bit).

\subsubsection{Pemisahan Kunci (Key Separation)}
Menggunakan kunci yang sama untuk operasi enkripsi dan autentikasi merupakan pelanggaran prinsip kriptografi yang dapat memicu interaksi antar-protokol tak terduga. Sistem menerapkan skema derivasi pemisahan kunci (\textit{Key Separation}) menggunakan master cipher sebagai fungsi pembangkit semu (\textit{pseudorandom function}/PRF):
\begin{align*}
K_{\text{enc}} &= \text{CIPH}_{K_{\text{master}}}(0^{15} \,\|\, \text{0x01}) \\
K_{\text{mac}} &= \text{CIPH}_{K_{\text{master}}}(0^{15} \,\|\, \text{0x02})
\end{align*}
Konstanta domain $0^{15} \,\|\, \text{0x01}$ dan $0^{15} \,\|\, \text{0x02}$ memastikan bahwa $K_{\text{enc}} \neq K_{\text{mac}}$ dengan independensi komputasional, sehingga kompromi pada satu operasi tidak menurunkan keamanan operasi lainnya.

\subsubsection{Perbandingan Waktu-Konstan (Constant-Time Comparison)}
Pengecekan kesamaan tag MAC rentan terhadap serangan kanal samping waktu (\textit{timing attack}) apabila menggunakan fungsi standar seperti \texttt{memcmp()} yang keluar lebih awal (\textit{early exit}) saat menemukan byte pertama yang berbeda. Sistem mengimplementasikan fungsi \texttt{util::constant\_time\_equal()}:
\[
\Delta = \bigvee_{i=0}^{15} (A_i \oplus B_i)
\]
Operasi akumulasi bitwise OR mengeksekusi iterasi penuh sepanjang 16 byte tanpa percabangan berbasis data (\textit{branchless}), sehingga waktu eksekusi identik baik untuk tag yang cocok maupun tag yang berbeda di byte pertama.

\subsection{Format Berkas Kemasan (Container v1)}

\subsubsection{Struktur Tata Letak Berkas}
Data terenkripsi dan metadata dibungkus dalam format biner terstandarisasi \textbf{Container v1} dengan tata letak deterministik sebagaimana tercantum pada Tabel~\ref{tab:container_layout}.

\begin{table}[ht]
\centering
\caption{Tata Letak Biner Kontainer Terenkripsi v1}
\label{tab:container_layout}
\begin{tabular}{|l|c|p{8.5cm}|}
\hline
\textbf{Field} & \textbf{Ukuran} & \textbf{Deskripsi} \\
\hline
Magic Bytes & 4 Byte & Pengenal format ASCII \texttt{"IF15"} (\texttt{0x49 0x46 0x31 0x35}) \\
Version & 1 Byte & Nomor versi protokol (\texttt{0x01}) \\
Mode ID & 1 Byte & Identifier mode cipher: \texttt{0x01}=ECB, \texttt{0x02}=CBC, \texttt{0x03}=CFB, \texttt{0x04}=OFB, \texttt{0x05}=CTR \\
Initialization Vector (IV) & 16 Byte & Entropi acak dari kernel OS (diisi nol pada mode ECB) \\
Ciphertext Payload & $N$ Byte & Data payload terenkripsi ($N \ge 0$, kelipatan 16 byte pada ECB/CBC) \\
CMAC Tag & 16 Byte & Tag autentikasi integritas RFC 4493 atas header dan ciphertext \\
\hline
\end{tabular}
\end{table}

Header kontainer memiliki ukuran tetap 22 byte, sehingga ukuran minimum berkas kontainer valid adalah $22 + 0 + 16 = 38$ byte (merepresentasikan plaintext kosong 0 byte pada mode stream CTR/CFB/OFB).

\subsubsection{Rationale Autentikasi Header dan IV}
Tag CMAC dihitung atas seluruh runtunan byte:
\[
\text{Data Terautentikasi} = \text{Magic (4B)} \,\|\, \text{Version (1B)} \,\|\, \text{Mode (1B)} \,\|\, \text{IV (16B)} \,\|\, \text{Ciphertext ($N$B)}
\]
Pengikatan header dan IV ke dalam lingkup autentikasi mencegah dua vektor serangan kritis:
\begin{enumerate}
    \item \textbf{Mode-Switching / Downgrade Attack}: Penyerang tidak dapat mengubah byte Mode ID (misalnya dari CTR ke ECB untuk membocorkan pola data) karena perubahan 1 bit pada byte mode akan menggagalkan verifikasi CMAC.
    \item \textbf{IV Substitution / Bit-Flipping Attack}: Pada mode CBC, pembalikan bit pada IV mentranslasikan pembalikan bit pada blok plaintext pertama hasil dekripsi. Dengan menyertakan IV dalam payload CMAC, setiap manipulasi IV ditolak sebelum dekripsi dimulai.
\end{enumerate}

\subsection{Antarmuka CLI dan Kebijakan Keamanan}

\subsubsection{Ringkasan Perintah CLI}
Aplikasi antarmuka baris perintah (\texttt{cipher\_cli}) menyediakan tiga subperintah operasional:
\begin{itemize}
    \item \texttt{keygen [-f|--force] -o <path>}: Membangkitkan kunci acak 128-bit (32 karakter heksadesimal).
    \item \texttt{encrypt [-f|--force] -m <mode> (-k <hex> | --key-file <path>) -i <in> -o <out>}: Mengenkripsi berkas masukan menjadi berkas kontainer terautentikasi.
    \item \texttt{decrypt [-f|--force] (-k <hex> | --key-file <path>) -i <in> -o <out>}: Memverifikasi integritas dan mendekripsi kontainer (mode dideteksi otomatis dari header kontainer).
\end{itemize}

\subsubsection{Tabel Kode Keluar (Exit Codes)}
Untuk memfasilitasi integrasi otomatisasi skrip dan pengujian sistem, \texttt{cipher\_cli} menetapkan konvensi kode keluar terstruktur seperti disajikan pada Tabel~\ref{tab:exit_codes}.

\begin{table}[ht]
\centering
\caption{Konvensi Kode Keluar (\textit{Exit Codes}) \texttt{cipher\_cli}}
\label{tab:exit_codes}
\begin{tabular}{|c|l|p{8.5cm}|}
\hline
\textbf{Kode} & \textbf{Kategori} & \textbf{Kondisi Pemicu} \\
\hline
0 & Sukses & Operasi enkripsi, dekripsi, atau keygen selesai dengan benar. \\
1 & Kesalahan CLI / Argumen & Opsi tidak dikenal, argumen wajib hilang, berkas target sudah ada tanpa \texttt{-f}, atau jalur \texttt{-i} sama dengan \texttt{-o}. \\
2 & Kesalahan I/O Berkas & Berkas masukan atau berkas kunci tidak ditemukan, gagal membaca atau menulis ke media penyimpanan. \\
3 & Format Rusak (\texttt{FormatError}) & Berkas bukan kontainer valid: magic salah, versi tidak didukung, mode tidak valid, panjang berkas $< 38$ byte. \\
4 & Integritas Gagal (\texttt{AuthError}) & Tag CMAC tidak valid akibat kunci salah atau berkas telah mengalami manipulasi/tampering. \\
\hline
\end{tabular}
\end{table}

\subsubsection{Kebijakan Keamanan Runtime}
Aplikasi menerapkan standar keamanan perangkat lunak kriptografi tingkat tinggi:
\begin{itemize}
    \item \textbf{Penghapusan Memori Kunci (\textit{Zeroization})}: Seluruh buffer yang menampung kunci mentah, representasi heksadesimal, string argumen baris perintah, dan plaintext segera dibersihkan menggunakan \texttt{util::secure\_zero()}. Fungsi ini memanfaatkan pointer \texttt{volatile} untuk mencegah optimasi kompilator (\textit{dead-store elimination}) menghapus pembersihan memori sensitif.
    \item \textbf{Perizinan Berkas POSIX 0600}: Pemanggilan \texttt{keygen} menetapkan mask berkas \texttt{S\_IRUSR | S\_IWUSR} melalui \texttt{open()} dan \texttt{fchmod()}, menjamin hanya pemilik proses yang memiliki izin membaca dan menulis berkas kunci.
    \item \textbf{Atomisitas Berkas dan Pencegahan Kebocoran}: Penulisan berkas keluaran dilakukan secara atomik melalui berkas sementara (\texttt{<path>.tmp}) sebelum di-rename. Apabila verifikasi integritas gagal (exit code 3 atau 4), aplikasi tidak pernah membuat atau meninggalkan berkas keluaran pada sistem berkas.
\end{itemize}

\subsection{Hasil Pengujian Sistem dan Evaluasi Kinerja}

Pengujian sistem dilakukan secara otomatis melalui rangkaian tes integrasi end-to-end (\texttt{tests/e2e.sh}) yang mencakup 59 pengujian fungsional dan integritas, serta pengujian unit pada utilitas (32 tes), CMAC (1011 tes), dan container (782 tes) dengan total 1884 assertions lulus 100\%.

\subsubsection{Uji Manipulasi Bit (1-Bit Tamper Resistance)}
Pengujian resistansi tampering dilakukan dengan membalik tepat satu bit (\textit{single-bit flip}) menggunakan operator XOR 1 pada berbagai posisi dalam kontainer terenkripsi pada seluruh lima mode cipher (ECB, CBC, CFB, OFB, CTR). Hasil pengujian empiris dirangkum pada Tabel~\ref{tab:tamper_results}.

\begin{table}[ht]
\centering
\caption{Hasil Uji Manipulasi 1-Bit pada Kontainer Terenkripsi}
\label{tab:tamper_results}
\begin{tabular}{|l|c|c|c|c|}
\hline
\textbf{Posisi Modifikasi} & \textbf{Offset Byte} & \textbf{Kode Keluar} & \textbf{Berkas Keluar Dibuat?} & \textbf{Tingkat Penolakan} \\
\hline
Magic Byte & 0 & 3 (\texttt{FormatError}) & Tidak Ada & 100\% (5/5 mode) \\
IV Byte & 6 & 4 (\texttt{AuthError}) & Tidak Ada & 100\% (5/5 mode) \\
Ciphertext Byte & 22 & 4 (\texttt{AuthError}) & Tidak Ada & 100\% (5/5 mode) \\
CMAC Tag Byte & Akhir & 4 (\texttt{AuthError}) & Tidak Ada & 100\% (5/5 mode) \\
Kunci Tidak Cocok & - & 4 (\texttt{AuthError}) & Tidak Ada & 100\% (5/5 mode) \\
\hline
\end{tabular}
\end{table}
Pada seluruh skenario di atas, aplikasi menolak operasi secara konsisten dan tidak meninggalkan berkas keluaran pada sistem berkas.

\subsubsection{Pengujian Batas Ukuran dan Berkas Asing}
Evaluasi keabsahan data dilakukan terhadap variasi ukuran payload:
\begin{itemize}
    \item \textbf{Berkas Kosong (0 Byte)}: Berhasil dienkripsi dan didekripsi kembali menghasilkan berkas 0 byte yang identik pada seluruh mode.
    \item \textbf{Ukuran Batas (Boundary Sizes)}: Ukuran 1, 15, 16, 17, 31, 32, dan 33 byte teruji menghasilkan kecocokan biner 100\% melalui perbandingan \texttt{cmp}.
    \item \textbf{Biner Arbitrer dan Skala Besar}: Runtunan 256 byte (\texttt{0x00} hingga \texttt{0xFF}) dan berkas acak berukuran 2 MiB berhasil diproses tanpa distorsi biner.
    \item \textbf{Berkas Terpotong (\textit{Truncated})}: Berkas kontainer yang dipotong 1 byte terakhir menghasilkan exit code 4, berkas dipotong hingga 37 byte ($< 38$ byte) menghasilkan exit code 3, dan berkas 0 byte menghasilkan exit code 3.
    \item \textbf{Berkas Asing}: Berkas acak 100 byte tanpa header \texttt{IF15} langsung ditolak dengan exit code 3.
\end{itemize}

\subsubsection{Pengukuran Kinerja Throughput}
Evaluasi performa throughput diukur pada payload acak berukuran 16 MiB (16.777.216 byte) pada arsitektur Apple Silicon (ARM64) dengan optimasi standar kompilator \texttt{-O3} setara. Hasil pengukuran disajikan pada Tabel~\ref{tab:perf_benchmark}.

\begin{table}[ht]
\centering
\caption{Kinerja Throughput Enkripsi dan Dekripsi (Payload 16 MiB)}
\label{tab:perf_benchmark}
\begin{tabular}{|l|c|c|c|c|}
\hline
\textbf{Mode Operasi} & \textbf{Waktu Enkripsi} & \textbf{Throughput Enkripsi} & \textbf{Waktu Dekripsi} & \textbf{Throughput Dekripsi} \\
\hline
CTR & 0,791 s & 20,22 MB/s & 0,791 s & 20,24 MB/s \\
CBC & 0,957 s & 16,72 MB/s & 0,794 s & 20,16 MB/s \\
\hline
\end{tabular}
\end{table}
Throughput round-trip mencapai 20,23 MB/s untuk mode CTR dan 18,28 MB/s untuk mode CBC. Perbedaan performa enkripsi CBC disebabkan oleh sifat sekuensial rantai blok CBC dan overhead penambahan padding PKCS\#7, sedangkan dekripsi CBC dan mode CTR dapat memproses data secara lebih efisien.

\subsection{Batasan Sistem dan Catatan Kriptografi}

\begin{enumerate}
    \item \textbf{Sifat Algoritma Blok Cipher Akademis}: Algoritma block cipher 128-bit yang digunakan merupakan desain khusus kurikulum perkuliahan. Algoritma ini belum melalui proses kompetisi kriptoanalisis publik bertahun-tahun (sebagaimana halnya AES/Rijndael), sehingga penggunaannya dibatasi untuk tujuan edukasi dan evaluasi akademik.
    \item \textbf{Kelemahan Pola Deterministik Mode ECB}: Pengujian sistem membuktikan secara empiris bahwa enkripsi dua blok 16-byte identik pada mode ECB menghasilkan dua blok ciphertext yang identik. Mode ECB tidak memberikan keamanan semantik (\textit{indistinguishability under chosen-plaintext attack}/IND-CPA). Pengguna sangat disarankan menggunakan mode CTR atau CBC dengan IV teracak untuk mengamankan data terstruktur.
    \item \textbf{Pemrosesan Berbasis Buffer Dalam Memori (\textit{In-Memory})}: Seluruh data masukan dimuat secara utuh ke dalam RAM (\texttt{std::vector<uint8\_t>}) sebelum enkripsi atau autentikasi dijalankan. Pendekatan ini optimal untuk berkas berukuran puluhan hingga ratusan megabyte, namun akan dibatasi oleh ketersediaan RAM fisik jika memproses berkas multi-gigabyte. Pengembangan di masa mendatang dapat mengintegrasikan arsitektur streaming berbasis \textit{chunked pipeline}.
\end{enumerate}
